%=====================================================================
\chapter{Resources and Further Reading}\label{app:resources}
%=====================================================================
\normalfigures

\section{The Papers This Book Is Built On}

Four of these are short enough to read in an evening and are the primary
sources for Part~I and Part~II. If you read nothing else from this appendix,
read the first.

\rowcolors{2}{white}{lightbg}
\begin{longtable}{@{}L{5.6cm}L{2.0cm}L{7.0cm}@{}}
\toprule
\rowcolor{primary!15}
\textbf{Source} & \textbf{Chapter} & \textbf{Why} \\
\midrule
\endhead
Popek \& Goldberg, \emph{Formal Requirements for Virtualizable Third Generation
Architectures}, CACM 17(7), 1974 & 3 & The theorem. Twelve pages, and the whole
of Part~II is a response to it \\
Robin \& Irvine, \emph{Analysis of the Intel Pentium's Ability to Support a
Secure Virtual Machine Monitor}, USENIX Security, 2000 & 3 & Where the
seventeen offending instructions were enumerated \\
Barham et al., \emph{Xen and the Art of Virtualization}, SOSP, 2003 & 2, 5 &
Paravirtualization, stated by the people who built it \\
Adams \& Agesen, \emph{A Comparison of Software and Hardware Techniques for
x86 Virtualization}, ASPLOS, 2006 & 3, 4 & Why binary translation beat the
first hardware assistance. The honest comparison, from VMware \\
Waldspurger, \emph{Memory Resource Management in VMware ESX Server}, OSDI,
2002 & 6 & Ballooning, page sharing and the reclamation ladder \\
Agache et al., \emph{Firecracker: Lightweight Virtualization for Serverless
Applications}, NSDI, 2020 & 11 & The microVM, and what was removed to get
125 ms \\
Clark et al., \emph{Live Migration of Virtual Machines}, NSDI, 2005 & 17 &
Pre-copy, and the convergence condition \\
Weil et al., \emph{Ceph: A Scalable, High-Performance Distributed File
System}, OSDI, 2006 & 12 & CRUSH, and placement without a lookup \\
Burns et al., \emph{Borg, Omega, and Kubernetes}, ACM Queue, 2016 & 10 & Why
Kubernetes is declarative, from the people who learned it the hard way \\
\bottomrule
\end{longtable}

\section{Specifications and Manuals}

The authoritative sources, for when a textbook account is not enough.

\begin{tight}
  \item \emph{Intel 64 and IA-32 Architectures Software Developer's Manual},
        Volume 3C --- VMX operation, the VMCS, EPT, VPID. Chapters 23--33.
  \item \emph{AMD64 Architecture Programmer's Manual}, Volume 2, Chapter 15 ---
        SVM, the VMCB, nested paging.
  \item \emph{Virtual I/O Device (VIRTIO) Specification}, OASIS --- the
        virtqueue, and every device type. \chref{paravirt}.
  \item \emph{Open Container Initiative} image, runtime and distribution
        specifications. \chref{docker}.
  \item \emph{PCI-SIG Single Root I/O Virtualization} specification ---
        physical and virtual functions. \chref{paravirt}.
  \item NIST SP 800-145, \emph{The NIST Definition of Cloud Computing} --- two
        pages, and the five characteristics of \chref{cloud}.
  \item NIST SP 800-125 series --- security guidance for full virtualization
        and for hypervisor deployment. \chref{security}.
\end{tight}

\section{Documentation Worth Reading Rather Than Searching}

\begin{tight}
  \item \textbf{libvirt} --- the domain XML reference. Every option in
        \chref{platforms}'s and \chref{paravirt}'s labs is documented there.
  \item \textbf{QEMU} --- the device model documentation, which doubles as a
        catalogue of the attack surface in \chref{security}.
  \item \textbf{Kubernetes} --- the concepts section, in order. It is unusually
        well written and it is the fastest route to \chref{kubernetes}.
  \item \textbf{Docker} --- the Dockerfile best-practices page. Short, and it is
        \chref{docker}'s build-cache rule stated by the people who implemented
        it.
  \item \textbf{VMware} --- the vSphere Resource Management guide, for
        \chref{cpumem}'s counters by their official names and thresholds.
  \item \textbf{Brendan Gregg's} systems-performance material, for the USE
        method of \chref{performance}.
\end{tight}

\section{Where to Practise}

\begin{tight}
  \item Your own laptop, with Appendix~A's stack. Everything in this book
        except three labs runs there.
  \item \textbf{Proxmox VE} on a spare machine --- a complete, free platform
        with clustering, backup and live migration. The closest thing to a
        vSphere experience without a licence.
  \item \textbf{kind} and \textbf{minikube} for Kubernetes.
  \item The three cloud free tiers of \chref{cloud}, with the budget alert of
        Appendix~A set first.
  \item \textbf{Katacoda-style browser labs} and the cloud providers' own
        guided labs --- useful for a first look, not a substitute for breaking
        something yourself.
\end{tight}

\section{Books}

\begin{tight}
  \item Smith \& Nair, \emph{Virtual Machines: Versatile Platforms for Systems
        and Processes}. The systematic treatment of the whole field, including
        the process-level virtual machines this course does not cover.
  \item Bugnion, Nieh \& Tsafrir, \emph{Hardware and Software Support for
        Virtualization}. Short, current, and the closest match to Parts~I
        and~II of this handout.
  \item Arpaci-Dusseau, \emph{Operating Systems: Three Easy Pieces}. Free
        online, and the prerequisite material for \chref{vmm} and
        \chref{containers} if those chapters were hard.
  \item Burns, Beda \& Hightower, \emph{Kubernetes: Up and Running}.
  \item Gregg, \emph{Systems Performance}. \chref{performance}'s method, at
        length.
\end{tight}

\section{On the Texts in the Published Outline}

\begin{alertbox}[Why Stanney and Burdea are not used]
The HEC 2017 outline for this course lists two set texts:

\begin{tight}
  \item K.\,M. Stanney (ed.), \emph{Handbook of Virtual Environments: Design,
        Implementation and Applications};
  \item G.\ Burdea \& P.\ Coiffet, \emph{Virtual Reality Technology}.
\end{tight}

Both are virtual-\emph{reality} references --- head-mounted displays, haptic
feedback, presence, simulator sickness. Neither discusses hypervisors,
virtual machine monitors, containers or any other topic in the outline's own
description, which is unambiguously about system virtualization.

\smallskip
The most likely explanation is that the reading list was assembled from the
course title rather than its contents. They are listed here for completeness
and because a student comparing this handout with the published outline
deserves to know why two named books are absent, not because they are useful
for this subject. If virtual reality is what you want, they are good books ---
for a different course.
\end{alertbox}

\section{Keeping Current}

Three chapters of this book date faster than the rest: \chref{cloud} (prices
and instance families), \chref{security} (vulnerabilities) and
\chref{confidential} (platform availability). \chref{gpu} is close behind.

\begin{tight}
  \item \textbf{Prices.} Always from the provider's own calculator, on the day.
        Every figure in \chref{cloud} is dated and says so.
  \item \textbf{Vulnerabilities.} The vendors' security advisory feeds --- VMware
        (Broadcom), Microsoft, Red Hat --- and CISA's Known Exploited
        Vulnerabilities catalogue, which is the one that tells you whether a
        flaw is being used rather than merely published.
  \item \textbf{Confidential computing.} The Confidential Computing Consortium,
        and each cloud's release notes; availability changed several times
        during 2026 alone.
  \item \textbf{Everything else.} The release notes of whatever you run. The
        mechanisms in Parts~I to~III have been stable for a decade and will
        outlast the products that implement them --- which is the reason this
        book is organised around mechanisms.
\end{tight}
